\documentclass[12pt]{article}
\usepackage[utf8]{inputenc}
\usepackage[T1]{fontenc}
\usepackage[margin=1in]{geometry}
\usepackage{setspace}
\usepackage{xcolor}
\usepackage[colorlinks=true, urlcolor=blue!60!black, linkcolor=blue!60!black]{hyperref}
\setlength{\parindent}{0em}
\setlength{\parskip}{0.6em}
\renewcommand{\baselinestretch}{1.0}
\begin{document}
\begin{center}
{\LARGE \textit{Nixon v. United States}, 506 U.S. 224 (1993)\par}
\vspace{0.5em}
{\large Case Brief --- Two-Page Edition\par}
\end{center}

\section{Memory Jogger}
The Senate's ``sole Power to try all Impeachments'' makes a removed judge's challenge to committee factfinding a nonjusticiable political question (\href{https://www.courtlistener.com/opinion/112803/nixon-v-united-states/}{CourtListener}).

\section{Detailed Case Facts}
Chief Judge Walter L. Nixon, Jr., convicted of lying to a grand jury, refused to resign. After the House impeached him, the Senate used Rule XI to have a committee take testimony and report; the full Senate heard final arguments and convicted him by more than two-thirds.

\section{Procedural History}
Nixon sued to void his conviction and restore his salary. The district court dismissed the case as a political question, and the D.C. Circuit affirmed.

\section{Judicial Votes}
Affirmed 9--0. Rehnquist wrote for six Justices; Stevens concurred; White (with Blackmun) and Souter concurred in the judgment.

\section{Holding}
The Impeachment Trial Clause textually commits trial procedure to the Senate, and ``try'' lacks judicially manageable standards, so a challenge to Rule XI is nonjusticiable. Express limits---oath, two-thirds, the Chief Justice presiding---may remain reviewable.

\section{Analysis of Opinions}
\textbf{Majority:} ``Try'' has broad meanings, the three express requirements imply no others, and ``sole'' means without interference. Impeachment is the Legislature's check on judges, who should not police it, and review would threaten finality. \textit{Powell v. McCormack} involved fixed textual limits. \textbf{White:} the claim is justiciable---otherwise the Senate could convict an official as ``a bad guy''---but committee factfinding is a valid trial. \textbf{Souter:} courts might act if the Senate convicted ``upon a coin toss.''

\section{Examples: Future Applications}
\textbf{Same-Side:} Senate limits on cross-examination are unreviewable. \textbf{Opposite-Side:} a conviction by less than two-thirds should be reviewable. \textbf{Fence-Sitter:} a summary conviction with no hearing at all.

\section{Critique}
Textualists praise respect for the Constitution's allocation of power; others say Rehnquist overread ``sole'' and turned ``try'' into a parchment barrier. Souter's carve-out requires the very merits judgment it disclaims.

\section{Key Quotations}
\begin{quote}
``The conclusion that the use of the word `try' in the first sentence of the Impeachment Trial Clause lacks sufficient precision to afford any judicially manageable standard of review of the Senate's actions is fortified by the existence of the three very specific requirements that the Constitution does impose on the Senate when trying impeachments: The Members must be under oath, a two-thirds vote is required to convict, and the Chief Justice presides when the President is tried.''
\end{quote}
\noindent --- Chief Justice Rehnquist, Opinion of the Court, 506 U.S. at 230.

\end{document}
